% sections/sec_feature.tex
%
% subfiles パッケージのサブファイル
% \documentclass[メインファイルへの相対パス]{subfiles} と書くことで
% このファイル単独でもコンパイル可能になる

\documentclass[../104_Subfiles.tex]{subfiles}

\begin{document}

\section{\texttt{subfiles} の特徴}

\texttt{subfiles} パッケージを使うと、各サブファイルが
\textbf{単独でもコンパイルできる}ようになります。

\begin{itemize}
  \item メインファイルからは \texttt{\textbackslash subfile\{...\}} で読み込む
  \item サブファイル単体でコンパイルすると、メインのプリアンブルが適用される
  \item \texttt{\textbackslash input} / \texttt{\textbackslash include} との最大の違いは
        \textbf{サブファイル単独コンパイルの可否}
\end{itemize}

\end{document}
